\subsection{LM25117 - U1}
\subsubsection{Timer Resistor \texorpdfstring{$R_\mathrm{T}$}{Rt}}
The switching frequency of the main buck regulator was set to \uoneFswSI\ via the $R_\mathrm{T}$ pin resistor; this enables sufficient separation from the adapter's lower switching frequency ($\approx \qty{100}{\kilo\hertz}$) to minimize beat frequency noise, while preventing excessive switching losses at higher frequencies.
\begin{align}
    R_\mathrm{T,calc}   &= \frac{5.2 \cdot 10^9}{f_\mathrm{sw}} - 949 \nonumber\\
                        &= \frac{5.2 \cdot 10^9}{\uoneFsw} - 949 \approx \uoneRtCalcSI \nonumber\\
    R_\mathrm{T,picked} &= \uoneRtSI
\end{align}

\subsubsection{Output Inductor \texorpdfstring{$L_\mathrm{o}$}{Lo}}
The output inductance was calculated using a nominal voltage of $V_\mathrm{out,typ} = \uoneVoutTypSI$ to achieve a balanced ripple current trade-off over the full output voltage range.
In addition, maximum inductor ripple ($I_\mathrm{pp}$) current was selected as \qty{20}{\percent} of the full load current.
\begin{align}
    L_\mathrm{o,calc}   &= \frac{V_\mathrm{out,typ}}{I_\mathrm{pp} \cdot f_\mathrm{sw}} \left(1 - \frac{V_\mathrm{out,typ}}{V_\mathrm{in}} \right) \nonumber\\
                        &= \frac{\uoneVoutTyp}{\uoneIpp \cdot \uoneFsw} \left(1 - \frac{\uoneVoutTyp}{\boardVin} \right) \approx \uoneLoCalcSI \nonumber\\
    L_\mathrm{o,picked} &= \uoneLoSI
\end{align}

\subsubsection{Current Sense Resistor \texorpdfstring{$R_\mathrm{CS}$}{Rcs}}
Maximum output current capability ($I_\mathrm{out,max}$) was selected as \qty{150}{\percent}. The $K$ factor defined in the datasheet was set to 1.
\begin{align}
    R_\mathrm{CS,calc}   &= \frac{V_\mathrm{CS,th}}{I_\mathrm{out,max} + \dfrac{V_\mathrm{out,typ} \cdot K}{f_\mathrm{sw} \cdot L_\mathrm{o}} - \dfrac{I_\mathrm{pp}}{2}} \nonumber\\
                         &= \frac{1.2}{\boardIoutMax + \dfrac{\uoneVoutTyp \cdot 1}{\uoneFsw \cdot \uoneLo} - \dfrac{\uoneIpp}{2}} \approx \uoneRcsCalcSI \nonumber\\
    R_\mathrm{CS,picked} &= \uoneRcsSI
\end{align}

\subsubsection{Ramp Resistor \texorpdfstring{$R_\mathrm{RAMP}$}{Rramp} and Ramp Capacitor \texorpdfstring{$C_\mathrm{RAMP}$}{Cramp}}
Ramp capacitor ($C_\mathrm{RAMP}$) was set to \uoneCrampSI.
\begin{align}
    R_\mathrm{RAMP,calc}   &= \frac{L_\mathrm{o}}{K \cdot C_\mathrm{RAMP} \cdot R_\mathrm{CS} \cdot A_\mathrm{RAMP}} \nonumber\\
                           &= \frac{\uoneLo}{1 \cdot \uoneCramp \cdot \uoneRcs \cdot 10} = \uoneRrampCalcSI \nonumber\\
    R_\mathrm{RAMP,picked} &= \uoneRrampSI
\end{align}

\subsubsection{UVLO Divider \texorpdfstring{$R_\mathrm{UV2}$}{Ruv2}, \texorpdfstring{$R_\mathrm{UV1}$}{Ruv1} and \texorpdfstring{$C_\mathrm{ft}$}{Cft}}
The desired startup voltage was set to \uoneVstartSI\ with \uoneVstartHysSI\ hysteresis to account for input adapter tolerances. 
Also, a \uoneCftSI\ capacitor was added in parallel to provide filtering for the divider.
\begin{align}
    R_\mathrm{UV2,calc}   &= \frac{V_\mathrm{start,hys}}{20 \cdot 10^{-6}} \nonumber\\
                          &= \frac{\uoneVstartHys}{20 \cdot 10^{-6}} = \uoneRuvCalctwoSI \nonumber\\
    R_\mathrm{UV2,picked} &= \uoneRuvtwoSI \\[1em]
    R_\mathrm{UV1,calc}   &= \frac{1.25 \cdot R_\mathrm{UV1}}{V_\mathrm{start} - 1.25} \nonumber\\
                          &= \frac{1.25 \cdot \uoneRuvone}{\uoneVstart - 1.25} \approx \uoneRuvCalconeSI \nonumber\\
    R_\mathrm{UV1,picked} &= \uoneRuvoneSI
\end{align}